\section{Developer setup and native packaging}\label{sec:build}
\subsection{Prepare a standalone checkout}
Use native 64-bit Python 3.12--3.13, JDK 17+ (\code{javac} and \code{jar}),
and optionally GNU Make. Mac packaging also needs Xcode command-line tools.
Windows requires x64 Python, including on an ARM host; verify
\code{sysconfig.get_platform()} reports \code{win-amd64}. No import or build
requires the former monorepo, old PyQt5 UI or an OpenRocket source checkout.

\begin{lstlisting}[language=bash]
# Run from RT_Telemetry_UI, not the report folder.
make setup PYTHON=python3.13
make run
make demo STATION=launch
make test
make lint
\end{lstlisting}
Setup creates \code{.venv}, installs the pinned build requirements and the
editable application. \code{uv sync --locked} is an alternative. Initial
downloads require internet; offline builds need staged matching wheels and
native vendor files. Metadata dependency ranges are not a tested interpreter
matrix. The tested package/runtime provenance is recorded below.

\subsection{Engine and private Java runtime}
\begin{lstlisting}[language=bash]
make runtime
make bridge OPENROCKET_DOWNLOAD=v6.2
# Alternative: a complete local team engine
make bridge OPENROCKET_JAR="/path/to/OpenRocket-MIT-v6.2.jar"
# Optional resolution at build time, not an app updater
make bridge OPENROCKET_DOWNLOAD=latest
\end{lstlisting}
Choose a local JAR or a download option, not both. The downloader queries
\code{MDNich/ActiveControl_MIT_RktTeam}, selects the exact
\code{OpenRocket-MIT-vX.Y.jar}, checks size and published SHA-256 when available,
and records provenance in \code{vendor/engine.json}. If no published digest
exists that fact is reported; a local hash is still recorded. Preparation
stages a candidate and compiles before replacing an existing engine/bridge.
Use a tested pinned release for production. A general upstream JAR may lack
the required fork fields and motor resources.

Runtime preparation verifies OS/architecture and downloads a checksum-verified
Temurin Java 17 runtime when needed. The bridge compiles for Java 17, records
engine/bridge hashes and retains source/provenance. Rebuild it after editing
\code{simulation_bridge/src/RocketBridge.java}. Never share \code{.venv} or
\code{vendor/java} across OS/architecture targets. The runtime and engine are
prepared at build time, not automatically downloaded by Settings or startup.

\subsection{Build with or without tests}
\begin{lstlisting}[language=bash]
# On macOS, after setup: tests + package + packaged verification
make -f Makefile.macos all OPENROCKET_DOWNLOAD=v6.2
# On Windows with GNU Make, after setup:
make -f Makefile.windows all OPENROCKET_DOWNLOAD=v6.2

# Package only; no test/verification targets:
make package-macos OPENROCKET_DOWNLOAD=v6.2
make package-windows OPENROCKET_DOWNLOAD=v6.2
\end{lstlisting}
Each command runs on its target OS; PyInstaller is not a cross-compiler.
An engine option stages the engine, bridge and matching runtime. Without an
option, packaging reuses existing vendor files: prepare runtime/bridge first.
The native \code{all} targets run test, package and verifier; run lint separately.
Each checkout's default PyInstaller cache is \code{build/pyinstaller-cache}.
An explicit \code{PYINSTALLER_CONFIG_DIR} overrides it; concurrent checkouts
must not share that writable cache.

Windows PowerShell can use the scripts without Make:
\begin{lstlisting}
py -3.13 scripts/build.py setup
.venv\Scripts\python.exe scripts/build.py runtime
.venv\Scripts\python.exe scripts/build.py bridge --download v6.2
.venv\Scripts\python.exe scripts/build.py package --platform Windows
# Optional validation, separate from the package-only command:
.venv\Scripts\python.exe -m pytest -q
.venv\Scripts\python.exe -m ruff check src tests --select F
.venv\Scripts\python.exe scripts/verify_package.py --model tests\fixtures\zephy_testlaunch.ork
\end{lstlisting}

\begin{longtable}{@{}L{.24\linewidth}L{.69\linewidth}@{}}
\toprule Target & Outputs under ignored \code{dist/} \\\midrule\endhead
Mac ARM64 & \code{RocketGNCMonitor-v0a.app}; \nolinkurl{RocketGNCMonitor-v0a-Darwin-arm64.dmg}. \\
Mac Intel & Same native app name in its separate checkout; \nolinkurl{RocketGNCMonitor-v0a-Darwin-x86_64.dmg}. \\
Windows x64 & \code{RocketGNCMonitor-v0a/} and \nolinkurl{RocketGNCMonitor-v0a-Windows-x64.zip}. \\
Universal Mac & \code{RocketGNCMonitor-v0a-universal.app}; \nolinkurl{RocketGNCMonitor-v0a-macOS-universal.dmg}; ZIP wrapping the DMG. \\
\bottomrule
\end{longtable}
Archives have SHA-256 sidecars. The bundle identity is
\code{edu.mit.rocketteam.gncmonitor.v0a}, version 0.1.1. The migration preserves
saved formats, application identity and separate v0a preferences. The PyInstaller
spec bundles Python/Qt, native MGRS, fonts, FFmpeg, Java, engine/bridge, resources
and notices, without unneeded WebEngine/QML or full scientific distributions.

\subsection{Universal Mac composition}
Build and verify ARM64 and x86\_64 apps in distinct native environments; Intel
can build/run through Rosetta on Apple Silicon. Then:
\begin{lstlisting}[language=bash]
make package-macos-universal \
  MAC_ARM64_APP="/path/to/arm64/RocketGNCMonitor-v0a.app" \
  MAC_X86_64_APP="/path/to/intel/RocketGNCMonitor-v0a.app"
\end{lstlisting}
\code{scripts/package_macos_universal.py} verifies architectures, signatures,
matching app versions/engines and Java provenance, preserves native inputs,
copies symlinks, compiles a small universal launcher and signs nested payloads
before the outer bundle. The launcher selects one complete native app and
uses \code{execv}, preserving arguments and exit status. This is not a claim
that every library is universal; concatenating finished PyInstaller archives
with \code{lipo} is not the build method. \code{UNIVERSAL_OUTPUT} changes the
destination. macOS 14+ reflects the included native runtime requirement.

\subsection{Verification, signing and cleanup}
Verify the extracted/relocated deliverable, including a path with spaces:
\begin{lstlisting}[language=bash]
.venv/bin/python scripts/verify_package.py \
  --package "/path with spaces/RocketGNCMonitor-v0a.app" \
  --model tests/fixtures/zephy_testlaunch.ork \
  --output build/package-verification
codesign --verify --deep --strict "/path/to/RocketGNCMonitor-v0a.app"
\end{lstlisting}
Use the Windows interpreter and extracted directory on Windows. The verifier
removes developer runtime overrides and uses minimal system PATH. Use native
desktop Qt for Windows; offscreen Mac regression tests do not replace native
visual checks. Universal verification must exercise both launcher slices.

Default Mac signing is ad-hoc. \code{MAC_SIGN_IDENTITY} selects an available
Developer ID identity; notarization is separate. Windows is unsigned unless a
separate signing workflow is applied. Preserve dependency licenses, Java legal
files, engine/bridge provenance and the user-designated proprietary font notice.
No signing credential belongs in source control.

Keep source, tests, small fixtures, authored documentation and provenance in
Git. Keep environments, vendor downloads, recordings, captures and distribution
binaries out. \code{make clean} in the \emph{application} removes both build
and dist, including finished packages. Remove only build when retaining final
deliverables. The \emph{report}'s separate Makefile clean removes its build
intermediates while keeping the final PDF.

\subsection{Manual Windows CI}
\src{.github/workflows/build-windows.yml} provides a manually dispatched Windows
Server 2025 build using x64 Python 3.13, JDK 17 and pinned requirements. It
prepares v6.2, runs tests/lint, packages and verifies extracted output with native
Qt/minimal PATH. Download the package/checksum artifact and verification artifact
promptly: retention is one day. The workflow does not publish a release or run
automatically. Recorded run metadata identifies the actual interpreter/JDK and
source revision; hosted runners do not validate physical USB or radio equipment.

\section{Testing and release evidence}\label{sec:tests}
\subsection{Regression responsibilities}
\begin{longtable}{@{}L{.31\linewidth}L{.62\linewidth}@{}}
\toprule Area & Required coverage \\\midrule\endhead
Protocol / transport & Golden command vectors, fragmentation/corruption, wrap/reboot, timeouts, uncertain writes, independent ports and polling. \\
Parity / controls & All original CSV columns and shortcuts, genuine Qt key dispatch, confirmations, no duplicate signal connections, role/uplink gating. \\
Station setup & Wizard cancellation, remember/forget/CLI precedence, immutable identity, Launch/Away boards, URRG presets, two-view shared ownership. \\
Video & Reservations through open/close, independent failure, station/channel assignments, promotion, stale source status, role isolation, local recording/replay. \\
Pose / graphics & Rejection of gyro integrals as attitude, quaternion/frame validation, unknown effects, reference provenance, upright camera, antenna chirality and interaction isolation. \\
Recording / flight & Queue/durability, active snapshots, finalized video, archive bounds/path/hash rejection, atomic failure, replay without original assets. \\
Coordinates / simulation & Datum and relative-location signs, MGRS/Plus Codes, wind, motor resolution, real example and Zephyrus engine runs, motion/events and virtual pointing. \\
Network / packaging & Bounded receive-only data, no remote command authority, runtime architecture, download integrity, failure preservation and relocated packaged execution. \\
\bottomrule
\end{longtable}
The captured legacy fixture records 118 rocket command cases across all
18 opcodes plus aliases, eight pointer cases, three tracking cases, 43 CSV
fields and ten shortcut strings, with original source hashes. It was captured
from the old class methods with physical connections replaced by memory streams.
Do not regenerate the fixture simply to accept a protocol regression.

Original conventions retained include five-decimal receiver coordinates,
zero receiver height for pointing, the receiver-height display division,
servo scales/rounding and good-CSV padding. Explicit fixes include rejected
packet counting, old bad-packet field inconsistencies, eight BMS columns and
duplicate shortcut connections. These fixtures require no old source checkout.

\subsection{0.1.1 validation baseline}
The following are completed migration/release checks from 25 October 2026,
recorded in \src{docs/MIGRATION_VALIDATION.md}. Building this report does not
rerun application/hardware qualification.
\begin{longtable}{@{}L{.30\linewidth}L{.63\linewidth}@{}}
\toprule Check & Recorded result \\\midrule\endhead
Mac Apple Silicon & 578 tests passed, no skips; regression suite used offscreen Qt. \\
Mac Intel / Rosetta & Full run: 573 passed, 5 skipped. Prepared-runtime simulation follow-up: all 5 simulation tests passed, covering four prior runtime skips. 577 Intel-applicable cases covered; the remaining ARM-host-only case passed on ARM. \\
Universal package & 20/20 packaged checks for ARM64 and 20/20 for Intel under Rosetta. Fresh path with spaces and minimal PATH. DMG payload audit: 1,743 files/symlinks matched; deep/strict signatures passed. \\
Windows clean runner & 567 tests passed, 11 platform skips; lint passed; 20/20 packaged checks. Windows Server 2025, Python 3.13.16 x64, Temurin 17.0.20.1+1. \\
Windows binary audit & ZIP checksum/CRC and embedded manifests matched; all 243 EXE/DLL/PYD binaries were AMD64. \\
Windows 11 ARM VM & Independent x64 source run: 567 passed, 11 skipped with Python 3.13.14 and v6.2. Adjusted virtual-pointer UI module: 7 passed. \\
\bottomrule
\end{longtable}
Packaged checks cover wizard/roles, camera inventories, locked LIVE startup,
all three demos, portable flights, URRG, virtual pointing, 3D ignition/recovery
and bundled-example/Zephyrus simulation. The successful Windows run is
\href{https://github.com/MIT-Rocket-Team/RT_Telemetry_UI/actions/runs/40213280209}{Actions run 40213280209}.
Its preceding failure exposed a UI-test visibility/timing issue; the test now
waits for visible/exposed controls before real mouse clicks, with behavior
assertions unchanged. Production application modules were unchanged.

The delivered Mac was built from \code{8edfdb0}; Windows from \code{b61a153}.
The release tag is \code{1923244}. Later commits contain workflow, test and
documentation changes, with identical production application code. Both
delivered archives use OpenRocket-MIT v6.2. A supplemental VM package using
v6.4 ultimately passed 20 checks after an unexplained transient video-start
failure and a separate guest-transport interruption; it is not the delivered ZIP.

\begin{longtable}{@{}L{.24\linewidth}L{.69\linewidth}@{}}
\toprule Release artifact & Integrity \\\midrule\endhead
Universal Mac DMG & 410,713,509 bytes. SHA-256:\newline\nolinkurl{acef899cf3b5eb2b43c402ce79c328df341f01db1d6c589c2eae45ad1eb862bc} \\
Windows x64 ZIP & 218,134,349 bytes. SHA-256:\newline\nolinkurl{5977838d86aaa9e32f57d717728c593f0e04713905d2b6dfc5b5ed912dba79f7} \\
\bottomrule
\end{longtable}
Intel was tested through Rosetta, not physical Intel hardware. Real cameras,
boards, pointer mechanics, APRS radio and LTU network were not connected.
Automated checks do not establish layout at every display scale, field
performance, notarization or qualified controlled-flight dynamics.

\section{Extension and maintenance guidance}\label{sec:extend}
\subsection{Hardware and new telemetry}
Specify framing, CRC, vehicle/source IDs, units, axes, timestamps/reset behavior
and channel mapping before adding an adapter. Populate explicit measured/demand
fields only when supplied. Preserve the Zephyrus adapter for old hardware/logs.
Add captured packets, malformed/reboot tests, source isolation, recording/archive
round trips and UI availability checks. Canard/tab names must not be assigned
speculatively to legacy channels. Individual-servo commands require a contract;
current encoders expose grouped roll/airbrake controls.

Define Launch enable and Iris selector commands, acknowledgements and timeout
semantics before replacing placeholders. Separate desired from confirmed state;
uncertain partial writes remain nonretrying. Measured pointer feedback needs a
structured parser distinct from last-sent target. Homing, cable limits and stop
must have firmware support rather than merely a renamed UI button.

\subsection{Data, simulation and presets}
Add presets to \code{site_presets.py} and the corresponding location UI, keeping
canonical coordinates, exact codes and explicit elevations. Test wizard defaults,
format changes, manual edits and persistence. Optional schema fields need defaults
and validation across mission/session/flight. Breaking semantics need a version
bump and migration/rejection, not silent unknown-field filtering.

Preserve archive limits, atomic replacement, active snapshot prefixes and
source provenance. Changes to the Java bridge require rebuilding against the
matching engine and recording options/input hashes. Controlled dynamics require
model qualification; do not clear \code{controlled_model_validated:false} merely
to suppress a label. Preserve raw inputs alongside any future fused estimate.

\subsection{Known follow-up work}
The release retains APRS, legacy terminology where it identifies the protocol,
and two Launch displays. Requested future changes include Launch-only YouTube
streaming, compact Windows layout improvements, a one-window/individual-servo
workflow, removal of APRS/legacy-facing labels, and a larger Away vertical plot.
These requests are not implemented by this documentation revision. Protocol
work for station transport, switches and sensor fusion remains separate.

For each change, check role/source boundaries and ownership, run meaningful
focused tests then release checks, inspect native windows at intended sizes
and representative smaller/high-DPI settings, rebuild the target package and
update operator/developer references. Never use physical recovery actuation
as a software smoke test.
